% Définition des couleurs
\usepackage{xcolor}

% Couleurs principales
\definecolor{primaryColor}{RGB}{0, 123, 255} % Bleu primaire
\definecolor{secondaryColor}{RGB}{108, 117, 125} % Gris secondaire
\definecolor{accentColor}{RGB}{255, 193, 7} % Jaune accent
\definecolor{successColor}{RGB}{40, 167, 69} % Vert succès
\definecolor{dangerColor}{RGB}{220, 53, 69} % Rouge danger

% Couleurs pour les chapitres et sections
\definecolor{chapterColor}{RGB}{0, 103, 185} % Bleu foncé pour les chapitres
\definecolor{sectionColor}{RGB}{51, 51, 153} % Bleu-violet pour les sections
\definecolor{subsectionColor}{RGB}{70, 130, 180} % Bleu ciel pour les sous-sections

% Couleurs pour les figures et tableaux
\definecolor{figBorder}{RGB}{200, 200, 200} % Gris clair pour les bordures de figures
\definecolor{figBackground}{RGB}{245, 245, 245} % Gris très clair pour l'arrière-plan des figures
\definecolor{captionColor}{RGB}{70, 70, 70} % Gris foncé pour les légendes

% Couleurs pour les UI éléments
\definecolor{uiFrameColor}{RGB}{230, 230, 230} % Gris clair pour les cadres UI
\definecolor{uiTitleColor}{RGB}{51, 51, 51} % Gris foncé pour les titres UI
